\bta{2021}



\begin{enumerate}
	%\renewcommand{\labelenumi}{\arabic{enumi}.}
	% A(\Alph) a(\alph) I(\Roman) i(\roman) 1(\arabic)
	%设定全局标号series=example	%引用全局变量resume=example
	%[topsep=-0.3em,parsep=-0.3em,itemsep=-0.3em,partopsep=-0.3em]
	%可使用leftmargin调整列表环境左边的空白长度 [leftmargin=0em]
	\item
	选择题 (本题满分 50 分，每小题 5 分)
\begin{enumerate}
	%\renewcommand{\labelenumi}{\arabic{enumi}.}
	% A(\Alph) a(\alph) I(\Roman) i(\roman) 1(\arabic)
	%设定全局标号series=example	%引用全局变量resume=example
	%[topsep=-0.3em,parsep=-0.3em,itemsep=-0.3em,partopsep=-0.3em]
	%可使用leftmargin调整列表环境左边的空白长度 [leftmargin=0em]
	\item
	$a=e^{x}, \quad b=1+x, \quad c=1+x+x^{2}$, 则在 $x=0$ 的$ \epsilon $无穷小邻域内, 大小估计正确的为 \xzanswer{A} 
	
\fourchoices
{$ b \leq a \leq c$}
{$a \leq b \leq c$}
{$b \leq c \leq a$}
{无法确定}

\item 
双曲正弦 $\sinh x=\frac{e^{x}-e^{-x}}{2} $，双曲余弦 $\cosh x=\frac{e^{x}+e^{-x}}{2}$，下列说法错误的是	 \xzanswer{B} 


\fourchoices
{双曲正弦为奇函数,双曲余弦为偶函数}
{双曲正弦为增函数,双曲余弦为减函数}
{双曲正弦的导为双曲余弦,双曲余弦的导为双曲正弦}
{$\cosh ^{2} x-\sinh ^{2} x=1$, 对任意 $x$ 总成立}

\item 
$n$ 为正整数, $\lim\limits _{n \rightarrow \infty} \cos 2 \pi \sqrt{n^{2}+n}=$ \xzanswer{C} 

\fourchoices
{$ 1 $}
{$ 0 $}
{$-1$}
{不存在}

\item 
设函数 $f(x)$ 是定义在 $(-1,1)$ 内的奇函数, 且 $\lim _{x \rightarrow 0^{+}} \frac{f(x)}{x}=a \neq 0$, 则 $f(x)$ 在 $x=0$ 处导数为 \xzanswer{A} 

\fourchoices
{$a$}
{$-a$}
{$ 0 $}
{不存在}

\item 
设向量 $\vec{a}$ 和 $\vec{b}$ 满足 $(\vec{a}+\vec{b}) \perp(\vec{a}-2 \vec{b}), \quad(3 \vec{a}-\vec{b}) \perp(\vec{a}+2 \vec{b})$, 则 $|\vec{a}|$ 与 $|\vec{b}|$ \xzanswer{D} 

\fourchoices
{$|\vec{a}|=\sqrt{\frac{1}{2}} |\vec{b}|$}
{$|\vec{a}|=\sqrt{2}|\vec{b}|$}
{$|\vec{a}|=\sqrt{\frac{2}{3}} |\vec{b}|$}
{$|\vec{a}|=\sqrt{\frac{3}{2}} |\vec{b}|$}

\item 
$I=\lim\limits _{x \rightarrow 0} \dfrac{\int_{\sin ^{3} x \cos x}^{e^{x^{2}}-1} \arctan \frac{3 t}{2+t} d t}{\arcsin x^{3}}$, 则 \xzanswer{D} 

\fourchoices
{$I$ 不存在}
{$I=3 / 2$}
{$I=1 / 2$}
{$I=0$}

\item 
$a_{n}>0,\left\{a_{n}\right\}$ 单调递减且趋于$  0  $, 则 $\sum\limits_{n=1}^{\infty}(-1)^{n-1} \sqrt{a_{n} \cdot a_{n-1}}$ \xzanswer{D} 
 

\fourchoices
{发散}
{绝对收敛}
{条件收敛}
{无法判断}

\item 
设函数 $z=x y+x F\left(\frac{y}{x}\right)$, 其中 $\mathrm{F}$ 为可导函数, 则 $x \frac{\partial z}{\partial x}+y \frac{\partial z}{\partial y}$ 的表达式是 \xzanswer{C} 

\fourchoices
{$Z-x y$}
{$ 0 $}
{$Z+x y$}
{$x y$}

\item 
$y^{\prime \prime}-6 y^{\prime}+8 y=e^{x}+e^{2 x}$ 的一个特解形式 \xzanswer{B} 


\fourchoices
{$a e^{x}+b e^{2 x} $}
{$ a e^{x}+b x e^{2 x} $}
{$a x e^{x}+b e^{2x} $}
{$ a x e^{x}+b x e^{2 x}$}

\item 
正确的是 \xzanswer{A、C} 


\fourchoices
{$\lim\limits _{s \rightarrow 0} \iint_{s <x^{2}+y^{2}< \frac{ 1 }{ 2 } } \frac{d x d y}{\left(x^{2}+y^{2}\right)\left(\ln \sqrt{x^{2}+y^{2}}\right)^{2}}$存在}
{$\lim\limits _{s \rightarrow 1} \iint_{-\frac{1}{2}<x^{2}+y^{2}<s} \frac{d x d y}{\left(x^{2}+y^{2}\right)\left(\ln \sqrt{x^{2}+y^{2}}\right)^{2}} $存在}
{$\lim\limits _{s \rightarrow 0} \iint_{s <x^{2}+y^{2}< \frac{ 1 }{ 2 } } \frac{\left(1+x^{2}\right) d x d y}{\left(x^{2}+y^{2}\right)\left(\ln \sqrt{x^{2}+y^{2}}\right)^{2}}$存在}
{以上均不对}




\end{enumerate}


\banswer{
\begin{enumerate}
	%\renewcommand{\labelenumi}{\arabic{enumi}.}
	\item
	\textbf{A}。由 $e^{x}=1+x+\frac{x^2}{2}+\frac{x^3}{6}+o\left(x^3\right)$ 得
	\[
	a-b=\frac{x^2}{2}+\frac{x^3}{6}+o\left(x^3\right)\geqslant0,
	\qquad
	c-a=\frac{x^2}{2}-\frac{x^3}{6}+o\left(x^3\right)\geqslant0
	\]
	（$|x|$ 充分小时 $x^2$ 项起主导作用，两者皆非负，$x=0$ 时取等号）。故 $b\leqslant a\leqslant c$。

	\item
	\textbf{B}。A 正确：$\sinh(-x)=-\sinh x$，$\cosh(-x)=\cosh x$。
	C 正确：$(\sinh x)'=\cosh x$，$(\cosh x)'=\sinh x$。
	D 正确：$\cosh^2x-\sinh^2x=\frac{\left(e^x+e^{-x}\right)^2-\left(e^x-e^{-x}\right)^2}{4}=1$。
	B 错误：$(\cosh x)'=\sinh x$ 在 $x<0$ 时为负、$x>0$ 时为正，
	故 $\cosh x$ 在 $(-\infty,0)$ 上递减、在 $(0,+\infty)$ 上递增，并非减函数。

	\item
	\textbf{C}。由
	$\sqrt{n^2+n}=n\sqrt{1+\frac1n}=n\left(1+\frac{1}{2n}-\frac{1}{8n^2}+o\left(\frac{1}{n^2}\right)\right)
	=n+\frac12-\frac{1}{8n}+o\left(\frac1n\right)$ 得
	\[
	2\pi\sqrt{n^2+n}=2n\pi+\pi-\frac{\pi}{4n}+o\left(\frac1n\right),
	\]
	故 $\cos2\pi\sqrt{n^2+n}=\cos\left(\pi-\frac{\pi}{4n}+o\left(\frac1n\right)\right)
	=-\cos\left(\frac{\pi}{4n}+o\left(\frac1n\right)\right)\to-1$。

	\item
	\textbf{A}。$f$ 为 $(-1,1)$ 上的奇函数故 $f(0)=0$。
	右导数 $\lim\limits_{x\to0^{+}}\frac{f(x)-f(0)}{x}=a$；
	左侧令 $x=-u\ (u\to0^{+})$，由奇性
	\[
	\frac{f(-u)-0}{-u}=\frac{-f(u)}{-u}=\frac{f(u)}{u}\to a .
	\]
	左右导数相等，故 $f'(0)=a$。

	\item
	\textbf{D}。两个垂直条件给出
	\begin{align*}
	(\vec{a}+\vec{b})\cdot(\vec{a}-2\vec{b})&=|\vec{a}|^2-\vec{a}\cdot\vec{b}-2|\vec{b}|^2=0,\\
	(3\vec{a}-\vec{b})\cdot(\vec{a}+2\vec{b})&=3|\vec{a}|^2+5\,\vec{a}\cdot\vec{b}-2|\vec{b}|^2=0 .
	\end{align*}
	由第一式 $\vec{a}\cdot\vec{b}=|\vec{a}|^2-2|\vec{b}|^2$，代入第二式得
	$8|\vec{a}|^2=12|\vec{b}|^2$，即 $|\vec{a}|=\sqrt{\frac{3}{2}}|\vec{b}|$。

	\item
	\textbf{D}。$t\to0$ 时 $\arctan\frac{3t}{2+t}\sim\frac{3t}{2}$。
	积分上、下限分别为 $U=e^{x^2}-1\sim x^2$ 与 $L=\sin^3x\cos x\sim x^3$，故
	\[
	\int_{L}^{U}\arctan\frac{3t}{2+t}\,dt\sim\int_{L}^{U}\frac{3t}{2}dt=\frac34\left(U^2-L^2\right)
	\sim\frac34x^4 ,
	\]
	而分母 $\arcsin x^3\sim x^3$，故 $I=\lim\limits_{x\to0}\frac{\frac34x^4}{x^3}=0$。

	\item
	\textbf{D（绝对或条件收敛类型不能确定）}。记 $b_n=\sqrt{a_na_{n-1}}$。由 $a_n$ 单调递减非负得
	$b_{n+1}=\sqrt{a_{n+1}a_n}\leqslant\sqrt{a_na_{n-1}}=b_n$，
	且 $0<b_n\leqslant a_{n-1}\to0$，故 $\left\{b_n\right\}$ 单调递减趋于 $0$。
	由 Leibniz 判别法，$\sum(-1)^{n-1}b_n$ 收敛。

	\par\vspace{1mm}\noindent
	\textbf{注：}{\kaishu
	严格地说，题设只能保证\emph{收敛}。是否绝对收敛并不确定：
	取 $a_n=\frac1{n+1}$（$n\ge0$）时 $b_n=\frac{1}{\sqrt{n(n+1)}}\sim\frac1n$，$\sum b_n$ 发散（条件收敛）；
	而取 $a_n=2^{-n}$ 时 $\sum b_n$ 收敛（绝对收敛）。
	因此不能断言条件收敛；D 仅表示无法确定收敛类型，并不表示无法判断是否收敛。
	}\par

	\item
	\textbf{C}。$xy$ 是 $2$ 次齐次函数，$xF\left(\frac{y}{x}\right)$ 是 $1$ 次齐次函数，
	由 Euler 齐次函数定理，
	\[
	x\frac{\partial z}{\partial x}+y\frac{\partial z}{\partial y}
	=2\cdot xy+1\cdot xF\left(\frac{y}{x}\right)=2xy+(z-xy)=z+xy .
	\]
	（直接求导亦可：$z_x=y+F-\frac{y}{x}F'$，$z_y=x+F'$，
	则 $xz_x+yz_y=xy+xF-yF'+xy+yF'=z+xy$。）

	\item
	\textbf{B}。特征方程 $r^2-6r+8=0$ 的根为 $r=2,4$。
	对 $e^{x}$：$1$ 不是特征根，特解取 $ae^{x}$；
	对 $e^{2x}$：$2$ 是单重特征根，需乘 $x$，特解取 $bxe^{2x}$。
	故形式为 $ae^{x}+bxe^{2x}$。

	\item
	\textbf{A、C}。用极坐标，$\frac{dxdy}{\left(x^2+y^2\right)\ln^2\sqrt{x^2+y^2}}
	=\frac{r\,drd\theta}{r^2\ln^2r}=\frac{drd\theta}{r\ln^2r}$，
	而由 $u=\ln r$ 得 $\int\frac{dr}{r\ln^2r}=-\frac{1}{\ln r}$。于是
	\[
	\iint\limits_{s<x^2+y^2<\frac12}=2\pi\left[-\frac{1}{\ln r}\right]_{\sqrt{s}}^{1/\sqrt2}
	=2\pi\left(\frac{1}{\ln\sqrt2}+\frac{1}{\ln\sqrt{s}}\right)
	\xrightarrow[s\to0^{+}]{}\frac{2\pi}{\ln\sqrt2}=\frac{4\pi}{\ln2},
	\]
	极限存在。而 B 中当 $s\to1$ 时 $-\frac{1}{\ln\sqrt{s}}\to\infty$，极限不存在。

	\par\vspace{1mm}\noindent
	\textbf{注：}{\kaishu
	按现有题面，选项 C 其实也收敛：多出的 $x^2=r^2\cos^2\theta$ 使被积函数
	增加一项 $\frac{r\cos^2\theta}{\ln^2r}$，它在 $r\to0^{+}$ 时趋于 $0$，可积。
	故按现有题面 A、C 同时成立，不能作为单选题唯一选 A；选项原文需核对。
	}\par
\end{enumerate}
}

	

\item 
(本题满分 10 分)
求常微分方程 $x y^{\prime \prime}-y^{\prime} \ln y^{\prime}+y^{\prime} \ln x=0$, 满足 $y(1)=2$ 和 $y^{\prime}(1)=e^{2}$ 的特解。

\banswer{
	设 $p=y^{\prime}$，原方程为 $xp^{\prime}=p\ln p-p\ln x=p\ln\frac{p}{x}$，即
	\[
	p^{\prime}=\frac{p}{x}\ln\frac{p}{x},
	\]
	这是齐次方程。令 $u=\frac{p}{x}$，则 $p=ux$，$p^{\prime}=u^{\prime}x+u$，代入得
	\[
	u^{\prime}x+u=u\ln u\quad\Longrightarrow\quad u^{\prime}x=u(\ln u-1).
	\]
	分离变量，再令 $v=\ln u-1$（$dv=\frac{du}{u}$）：
	\[
	\int\frac{dv}{v}=\int\frac{dx}{x}\quad\Longrightarrow\quad \ln u-1=C_1x ,
	\]
	即 $\frac{p}{x}=e^{1+C_1x}$，$y^{\prime}=p=xe^{1+C_1x}$。

	由 $y^{\prime}(1)=e^{2}$ 得 $e^{1+C_1}=e^{2}$，故 $C_1=1$，$y^{\prime}=xe^{1+x}$。积分：
	\[
	y=\int xe^{1+x}dx=e\int xe^{x}dx=e(x-1)e^{x}+C_2=(x-1)e^{x+1}+C_2 .
	\]
	由 $y(1)=2$ 得 $C_2=2$。所以
	\[
	y=(x-1) e^{x+1}+2 .
	\]
}



\item 
(本题满分 10 分)
求过直线 $l_{1}: \frac{x-1}{2}=\frac{y}{1}=\frac{z+1}{3}$ 且平行于 $l_{2}:\left\{\begin{array}{l}2 x+y-z+1=0 \\ x-2 y+z-2=0\end{array}\right.$ 的平面方程。

\banswer{
	所求平面过 $l_1$，故含 $l_1$ 的方向 $\bm{s}_1=(2,1,3)$ 与点 $M_0(1,0,-1)$；
	又平行于 $l_2$，故还平行于 $l_2$ 的方向
	\[
	\bm{s}_2=(2,1,-1)\times(1,-2,1)=(-1,-3,-5) .
	\]
	于是平面法向量为
	\[
	\bm{n}=\bm{s}_1\times\bm{s}_2=
	\begin{vmatrix}
	\bm{i} & \bm{j} & \bm{k}\\
	2 & 1 & 3\\
	-1 & -3 & -5
	\end{vmatrix}
	=(4,\ 7,\ -5).
	\]
	由点法式得
	\[
	4(x-1)+7 y-5(z+1)=0 ,\qquad\text{即}\quad 4x+7y-5z-9=0 .
	\]
	（校验：$M_0$ 满足方程；且 $\bm{n}\cdot\bm{s}_2=-4-21+25=0$，故平面确与 $l_2$ 平行。）
}



\item 
(本题满分 10 分)
设 $f(x)=\lim _{n \rightarrow+\infty} \frac{x^{2 n-1}+n x \sin \frac{x}{n}}{x^{2 n}+1}$, 讨论 $f(x)$ 连续性。

\banswer{
	先求出 $f$ 的显式表达式。注意 $n\to+\infty$ 时
	$n x\sin\frac{x}{n}=x\cdot\frac{\sin(x/n)}{1/n}\to x^{2}$，再按 $|x|$ 与 $1$ 的大小分段：
	\[
	f(x)=
	\begin{cases}
	\dfrac{0+x^{2}}{0+1}=x^{2}, & |x|<1 \quad\left(x^{2n}\to0\right),\\[8pt]
	\dfrac{x^{-1}+0}{1+0}=\dfrac{1}{x}, & |x|>1 \quad(\text{分子分母同除 }x^{2n}),\\[8pt]
	\dfrac{x^{2n-1}+x^{2}}{2}, & |x|=1 .
	\end{cases}
	\]
	当 $|x|=1$ 时 $x^{2n}=1$、$x^{2n-1}=\frac1x=x$，故 $f(1)=\frac{1+1}{2}=1$，
	$f(-1)=\frac{-1+1}{2}=0$。

	\textbf{逐点讨论.}
	\begin{itemize}
		\item 在 $(-1,1)$ 内 $f(x)=x^{2}$ 连续；在 $|x|>1$ 上 $f(x)=\frac1x$ 连续。
		\item $x=1$：左极限 $\lim\limits_{x\to1^{-}}x^{2}=1$，右极限 $\lim\limits_{x\to1^{+}}\frac1x=1$，
		而 $f(1)=1$，三者相等，故 $f$ 在 $x=1$ \textbf{连续}。
		\item $x=-1$：左极限 $\lim\limits_{x\to-1^{-}}\frac1x=-1$，
		右极限 $\lim\limits_{x\to-1^{+}}x^{2}=1$，而 $f(-1)=0$。
		三者互不相等，故 $x=-1$ 是 \textbf{跳跃间断点}（第一类间断点）。
	\end{itemize}

	\textbf{结论.} $f(x)$ 在 $\mathbb{R}\setminus\{-1\}$ 上处处连续，
	在 $x=-1$ 处间断（跳跃间断点）。

	\par\vspace{1mm}\noindent
	\textbf{注：}{\kaishu
	本题原答案称“处处连续”，有误——漏掉了 $x=-1$。
	关键在于 $x^{2n-1}$ 在 $x=-1$ 处为 $-1$ 而非 $+1$，使 $f(-1)=0$
	既不等于左极限 $-1$，也不等于右极限 $+1$。
	}\par
}


\item 
(本题满分 10 分)
设 $g(x)=\left\{\begin{array}{cc}\frac{e^{x}-1}{x} & x \neq 0 \\ 1 & x=0\end{array} \right.$,$f(x)=g^{\prime}(x)$。求 $f^{(n)}(0)$。

\banswer{
	用幂级数展开最省事。由 $e^{x}-1=\sum\limits_{k=1}^{\infty}\frac{x^{k}}{k!}$ 得
	\[
	g(x)=\frac{e^{x}-1}{x}=\sum_{k=1}^{\infty}\frac{x^{k-1}}{k!}=\sum_{m=0}^{\infty}\frac{x^{m}}{(m+1)!},
	\]
	该式在 $x=0$ 处取值 $1$，与 $g(0)=1$ 一致，故它在整个 $\mathbb{R}$ 上成立（且 $g$ 解析）。
	逐项求导，
	\[
	f(x)=g^{\prime}(x)=\sum_{m=1}^{\infty}\frac{m\,x^{m-1}}{(m+1)!}
	=\sum_{k=0}^{\infty}\frac{(k+1)\,x^{k}}{(k+2)!} .
	\]
	幂级数的系数即 Taylor 系数，故 $\frac{f^{(n)}(0)}{n!}=\frac{n+1}{(n+2)!}$，于是
	\[
	f^{(n)}(0)=\frac{n!\,(n+1)}{(n+2)!}=\frac{n!\,(n+1)}{(n+2)(n+1)\,n!}=\frac{1}{n+2} .
	\]
	（例如 $f(0)=\frac12$，$f^{\prime}(0)=\frac13$，与直接计算一致。）
}


\newpage
\item 
(本题满分 10 分)
设曲面 $\Sigma=\left\{(x, y, z) \in R^{3} \mid x^{2}+y^{2}+z^{2}=1\right\}$ 的外侧, 计算积分
\[ 
\iint_{\Sigma} y^{2} z d x d y+x^{2} y d z d x
 \]

\banswer{
	把积分写成标准形式 $\iint_\Sigma P\,dydz+Q\,dzdx+R\,dxdy$，则
	\[
	P=0,\qquad Q=x^{2}y,\qquad R=y^{2}z .
	\]
	$\Sigma$ 是闭曲面且取外侧，故可直接用高斯公式。
	\[
	\operatorname{div}\bm{F}=\frac{\partial P}{\partial x}+\frac{\partial Q}{\partial y}+\frac{\partial R}{\partial z}
	=0+x^{2}+y^{2} .
	\]
	设 $\Omega$ 为单位球体，则
	\[
	I=\iiint\limits_{\Omega}\left(x^{2}+y^{2}\right)dV .
	\]
	由对称性 $\iiint_\Omega x^2dV=\iiint_\Omega y^2dV=\iiint_\Omega z^2dV
	=\frac13\iiint_\Omega\left(x^2+y^2+z^2\right)dV$，而用球坐标
	\[
	\iiint\limits_{\Omega}\left(x^{2}+y^{2}+z^{2}\right)dV=\int_0^1 r^{2}\cdot4\pi r^{2}dr=\frac{4\pi}{5},
	\]
	故 $\iiint_\Omega x^2dV=\frac{4\pi}{15}$，从而
	\[
	I=2\times\frac{4\pi}{15}=\frac{8\pi}{15} .
	\]

	\par\vspace{1mm}\noindent
	\textbf{注：}{\kaishu
	本题原答案作 $\frac{3}{5}\pi=\frac{9\pi}{15}$，有误；正确值为
	$\frac{8\pi}{15}\approx1.6755$，数值积分亦予证实。
	}\par
}



\item 
(本题满分 10 分)
$\lim _{x \rightarrow 0} \frac{a x+\sin x}{\int_{b}^{x} \frac{\ln \left(1+t^{3}\right)}{t} d t}=c,(c \neq 0)$, 求出 $a, b, c$ 的值。

\banswer{
	记分母 $\Phi(x)=\int_{b}^{x}\frac{\ln\left(1+t^{3}\right)}{t}dt$，被积函数在 $t\to0$ 时
	$\frac{\ln\left(1+t^{3}\right)}{t}\sim\frac{t^{3}}{t}=t^{2}$。

	\textbf{定 $b$.} 分子 $ax+\sin x\to0\ (x\to0)$。若 $b\neq0$，则
	$\Phi(0)=\int_b^0\frac{\ln(1+t^3)}{t}dt\neq0$（被积函数在 $0$ 与 $b$ 之间不变号），
	于是极限为 $\frac{0}{\Phi(0)}=0=c$，与 $c\neq0$ 矛盾。故必须 $b=0$，此时
	\[
	\Phi(x)=\int_{0}^{x}\frac{\ln\left(1+t^{3}\right)}{t}dt\sim\int_0^x t^{2}dt=\frac{x^{3}}{3}.
	\]

	\textbf{定 $a$.} 分子展开 $ax+\sin x=(a+1)x-\frac{x^{3}}{6}+o\left(x^{3}\right)$。
	分母是 $x^{3}$ 阶，若 $a+1\neq0$ 则分子是 $x$ 阶，商趋于 $\infty$，极限不存在。
	故 $a+1=0$，即 $a=-1$。

	\textbf{定 $c$.} 此时
	\[
	c=\lim_{x\to0}\frac{-\frac{x^{3}}{6}+o\left(x^{3}\right)}{\frac{x^{3}}{3}+o\left(x^{3}\right)}
	=-\frac{1/6}{1/3}=-\frac{1}{2} .
	\]
	所以 $a=-1, b=0, c=-\frac{1}{2}$。
}



\item 
(本题满分 10 分)
设二元函数 $f(x, y)=\left\{\begin{array}{ll}\frac{\sqrt{|x y|}}{x^{2}+y^{2}} \sin \left(x^{2}+y^{2}\right) & x^{2}+y^{2} \neq 0 \\ 0 & x^{2}+y^{2}=0\end{array}\right.$，讨论$ f(x,y) $在$ (0,0) $处的可微性。	

\banswer{
	记 $\rho=\sqrt{x^{2}+y^{2}}$。

	\textbf{先求两个偏导数.} 沿坐标轴 $xy=0$，故 $f(x,0)\equiv0$、$f(0,y)\equiv0$，于是
	\[
	f_x(0,0)=\lim_{x\to0}\frac{f(x,0)-0}{x}=0,\qquad f_y(0,0)=0 .
	\]

	\textbf{再验证增量.} 若 $f$ 在 $(0,0)$ 可微，则必有
	$\frac{f(x,y)-f(0,0)-f_x(0,0)x-f_y(0,0)y}{\rho}=\frac{f(x,y)}{\rho}\to0$。而
	\[
	\frac{f(x,y)}{\rho}=\frac{\sqrt{|xy|}}{\rho}\cdot\frac{\sin\rho^{2}}{\rho^{2}} ,
	\]
	其中 $\frac{\sin\rho^{2}}{\rho^{2}}\to1$。沿 $y=x\ (x\to0^{+})$，
	\[
	\frac{\sqrt{|xy|}}{\rho}=\frac{|x|}{\sqrt{2}\,|x|}=\frac{1}{\sqrt{2}} ,
	\]
	故 $\frac{f(x,y)}{\rho}\to\frac{1}{\sqrt2}\neq0$；而沿 $y=0$ 该比值恒为 $0$。
	极限不存在（更不为 $0$），故增量不是 $o(\rho)$。

	\textbf{结论.} $f(x, y)$ 在 $(0,0)$ 点处\textbf{不可微}（尽管两个偏导数都存在）。
}




\item 
(本题满分 10 分)
$f(x)=\left\{\begin{array}{cc}x & -\pi \leq x < 0 \\ 0 & 0 \leq x \leq \pi\end{array}\right.$，将其展开为傅里叶级数。


\banswer{
	按周期 $2\pi$ 展开，需要 $a_0,a_n,b_n$。因 $f$ 在 $[0,\pi]$ 上为零，各积分只剩 $[-\pi,0]$。

	\textbf{常数项.}
	\[
	a_0=\frac{1}{\pi}\int_{-\pi}^{\pi}f(x)dx=\frac{1}{\pi}\int_{-\pi}^{0}x\,dx
	=\frac{1}{\pi}\left[\frac{x^{2}}{2}\right]_{-\pi}^{0}=-\frac{\pi}{2},
	\qquad\text{故常数项}\ \frac{a_0}{2}=-\frac{\pi}{4}.
	\]

	\textbf{余弦系数.} 由 $\int x\cos nx\,dx=\frac{x\sin nx}{n}+\frac{\cos nx}{n^{2}}$，
	\[
	a_n=\frac{1}{\pi}\int_{-\pi}^{0}x\cos nx\,dx
	=\frac{1}{\pi}\left(\frac{1}{n^{2}}-\frac{(-1)^{n}}{n^{2}}\right)
	=\frac{1-(-1)^{n}}{n^{2}\pi}
	=\begin{cases}\dfrac{2}{n^{2}\pi}, & n\ \text{为奇数},\\[4pt]0,& n\ \text{为偶数}.\end{cases}
	\]

	\textbf{正弦系数.} 由 $\int x\sin nx\,dx=-\frac{x\cos nx}{n}+\frac{\sin nx}{n^{2}}$，
	\[
	b_n=\frac{1}{\pi}\int_{-\pi}^{0}x\sin nx\,dx
	=\frac{1}{\pi}\left(0-\frac{\pi(-1)^{n}}{n}\right)=\frac{(-1)^{n+1}}{n} .
	\]

	\textbf{结论.}
	\[
	f(x)=-\frac{\pi}{4}+\sum\limits_{n=1}^{\infty}\left\{\frac{\left[1-(-1)^{n}\right]}{n^{2} \pi} \cos nx+\frac{(-1)^{n+1}}{n} \sin nx\right\},
	\]
	在 $f$ 的连续点处成立；特别地 $x=0$ 是连续点，级数值为 $f(0)=0$。
	周期延拓在 $x=\pm\pi$ 处有跳跃，级数收敛到左右极限的平均值 $-\frac{\pi}{2}$。

	\par\vspace{1mm}\noindent
	\textbf{注：}{\kaishu
	原答案常数项作 $+\frac{\pi}{4}$，符号有误（应为 $-\frac{\pi}{4}$，因 $a_0=-\frac{\pi}{2}$）；
	另外求和号内的 $\cos x$、$\sin x$ 应为 $\cos nx$、$\sin nx$。
	取 $x=0$ 可以校验：级数值 $-\frac{\pi}{4}+\sum\limits_{n\ \text{奇}}\frac{2}{n^{2}\pi}
	=-\frac{\pi}{4}+\frac{2}{\pi}\cdot\frac{\pi^{2}}{8}=0$，
	正是 $f(0^{-})$ 与 $f(0^{+})$ 的平均值。
	}\par
}


\item 
(本题满分 10 分)
设函数$ f(x) $在$ [a,b] $上连续，在$ (a,b) $内可导，证明：如果$ f(x) $为非线性函数，
则存在 $\exists \ \xi_{1}, \xi_{2} \in(a, b)$ 使$f^{\prime}(\xi_{1})=\frac{f(b)-f(a)}{b-a}$，$f^{\prime}\left(\xi_{2}\right)<\frac{f(b)-f(a)}{b-a}$。

\banswer{
	记 $k=\frac{f(b)-f(a)}{b-a}$。由 Lagrange 中值定理，$\exists\,\xi_{1}\in(a,b)$ 使得 $f'(\xi_{1})=k$。\\
	令 $g(x)=f(x)-[f(a)+k(x-a)]$，则 $g\in C[a,b]\cap D(a,b)$，$g(a)=g(b)=0$ 且 $g\not\equiv 0$（$f$ 非线性）。
	若 $\exists\,x_{0}\in(a,b)$ 使 $g(x_{0})>0$，由连续函数最值定理，$g$ 在 $(a,b)$ 内某点 $c$ 取得正最大值，故 $g'(c)=0$。在 $[c,b]$ 上对 $g$ 用 Lagrange 定理：
	$\exists\,\xi_{2}\in(c,b)$ 使 $g'(\xi_{2})=\frac{g(b)-g(c)}{b-c}=-\frac{g(c)}{b-c}<0$，即 $f'(\xi_{2})<k$。\\
	若 $\exists\,x_{0}\in(a,b)$ 使 $g(x_{0})<0$，则 $g$ 在 $(a,b)$ 内某点 $c$ 取得负最小值，$g'(c)=0$。在 $[a,c]$ 上对 $g$ 用 Lagrange 定理：
	$\exists\,\xi_{2}\in(a,c)$ 使 $g'(\xi_{2})=\frac{g(c)-g(a)}{c-a}=\frac{g(c)}{c-a}<0$，亦有 $f'(\xi_{2})<k$。
	综上得证。
}



\item 
(本题满分 10 分)
已知函数 $z=f(x, y)$ 的全微分 $d z=2 x d x-2 y d y$ 并且 $f(1,1)=2020$; 求出
$f(x, y)$ 在椭圆域 $D=\left\{(x, y) \mid x^{2}+\frac{y^{2}}{4} \leq 1\right\}$ 上最大值和最小值。

\banswer{
	\textbf{先求 $f$.} 由 $dz=2xdx-2ydy$ 得 $f_x=2x$、$f_y=-2y$，故
	$f(x,y)=x^{2}-y^{2}+C$。由 $f(1,1)=1-1+C=2020$ 得 $C=2020$，即
	\[
	f(x,y)=x^{2}-y^{2}+2020 .
	\]

	\textbf{内部驻点.} $\nabla f=(2x,-2y)=\bm{0}$ 只有 $(0,0)$，$f(0,0)=2020$；
	因 $f_{xx}f_{yy}-f_{xy}^{2}=2\cdot(-2)-0=-4<0$，这是鞍点，不是极值点。
	故最值必在边界取得。

	\textbf{边界.} 在 $x^{2}+\frac{y^{2}}{4}=1$ 上有 $y^{2}=4\left(1-x^{2}\right)$，代入得
	\[
	f=x^{2}-4\left(1-x^{2}\right)+2020=5x^{2}+2016,\qquad x\in[-1,1] .
	\]
	这是关于 $x^{2}$ 的增函数，于是
	\[
	\max f=5+2016=2021\ \ \text{在}\ x=\pm1,\ y=0\ \text{处取到},
	\]
	\[
	\min f=2016\ \ \text{在}\ x=0,\ y=\pm2\ \text{处取到}.
	\]

	\textbf{结论.} 最大值 $f(\pm 1,0)=2021$；最小值 $f(0, \pm 2)=2016$
	（内部鞍点值 $2020$ 介于两者之间，不影响最值）。
}



	
	
	
\end{enumerate}


